% !TeX program = pdflatex
% Team 14 — MPO Zadanie 2
% report_tim14_2.tex  (cielovy rozsah 4–6 stran)
\documentclass[11pt,a4paper]{article}

\usepackage[T1]{fontenc}
\usepackage[utf8]{inputenc}
\usepackage[slovak]{babel}
\usepackage{lmodern}
\usepackage{amsmath,amssymb}
\usepackage{booktabs}
\usepackage{tabularx}
\usepackage{geometry}
\usepackage{setspace}
\usepackage{enumitem}
\usepackage{microtype}
\usepackage{hyperref}
\usepackage{fancyhdr}
\usepackage{titlesec}

\geometry{margin=2.2cm}
\setstretch{1.08}
\setlength{\parskip}{0.35em}
\setlength{\parindent}{0pt}
\setlist[itemize]{leftmargin=1.4em,itemsep=0.15em,topsep=0.25em}
\setlist[enumerate]{leftmargin=1.6em,itemsep=0.15em,topsep=0.25em}
\titleformat{\section}{\normalsize\bfseries}{\thesection.}{0.5em}{}
\titleformat{\subsection}{\normalsize\bfseries}{\thesubsection}{0.5em}{}
\titlespacing*{\section}{0pt}{0.9em}{0.35em}
\titlespacing*{\subsection}{0pt}{0.6em}{0.2em}

\pagestyle{fancy}
\fancyhf{}
\fancyhead[L]{\small FEI KKUI}
\fancyhead[R]{\small Tim~14 --- Zadanie~2}
\fancyfoot[C]{\thepage}
\renewcommand{\headrulewidth}{0.3pt}

\hypersetup{
  colorlinks=true,
  linkcolor=black,
  urlcolor=black,
  citecolor=black,
  pdftitle={MPO Zadanie 2 — Tim 14},
  pdfauthor={Tim 14}
}

\begin{document}

\begin{center}
  {\large\bfseries Moderné prístupy v~optimalizácii}\\[0.35em]
  {\bfseries Zadanie~2: Statická optimalizácia vlakovej trasy\\
  s~penalizáciou prestupov}\\[0.55em]
  {\small Tim~14 \quad|\quad FEI KKUI TUKE \quad|\quad November 2026}
\end{center}
\vspace{0.4em}
\hrule
\vspace{0.7em}

%==============================================================================
\section{Formulácia problému a~revízie podľa spätnej väzby}
%==============================================================================

Cieľom je nájsť najlepšiu prípustnú vlakovú trasu medzi východiskovou stanicou~$s$
a~cieľovou stanicou~$z$ pre zvolený deň~$D$ a~najskorší prípustný odchod~$t_0$.
Trasa môže byť priama alebo môže obsahovať viacero vlakov s~prestupmi.
Pracujeme so \emph{statickým} cestovným poriadkom~$J$ (bez meškaní a~rušení).
Prvý odchod musí spadať do okna~$[t_0,\,t_0+H]$ (predvolene $H=120$\,min);
čakanie pred prvým vlakom sa nehodnotí, čakanie pri prestupoch áno.
Minimálny prestupný interval je $\tau_{\min}=5$\,min; návrat do už navštívenej
stanice nepovoľujeme.

Kandidátne riešenie je usporiadaná postupnosť~$k$ priamych úsekov (legov)
\begin{equation}
  x=(\ell_1,\ell_2,\ldots,\ell_k),\qquad
  k\ge 1,\qquad
  \ell_i=(id_i,\,u_i,\,v_i,\,d_i,\,a_i),
\end{equation}
kde $id_i$ je identifikátor vlaku, $u_i$/$v_i$ nástupná/výstupná stanica
a~$d_i$/$a_i$ časy odchodu/príchodu v~celých minútach od polnoci dňa~$D$
(napr.\ $08{:}10\mapsto 490$, nasledujúci deň $00{:}30\mapsto 1470$).
Účelová funkcia zostáva jednoúčelová:
\begin{equation}
  T(x)=a_k-d_1,\qquad
  N(x)=k-1,\qquad
  f(x)=T(x)+\lambda\cdot N(x)\to\min,
\end{equation}
s~predvoleným $\lambda=15$\,min. Pri rovnakom~$f(x)$ uprednostníme skorší
konečný príchod~$a_k$ a~následne skorší prvý odchod~$d_1$, t.\ j.\ porovnávame
trojicu $(f(x),\,a_k,\,d_1)$.

\paragraph{Revízie podľa pripomienok k~Zadaniu~1.}
\begin{enumerate}[label=\arabic*)]
  \item \textbf{Zdroj cestovného poriadku.}
        Používame filtrovaný koridorový podmnožinu v~štýle ZSSK~GTFS
        (Košice -- Prešov -- Poprad-Tatry -- Vrútky -- Žilina -- Trenčín --
        Trnava -- Bratislava) a~doplňujúci deterministický syntetický dataset
        na izoláciu hraničných prípadov (prechod cez polnoc, tesný prestup,
        nedosiahnuteľný cieľ).

  \item \textbf{Ploché indexované štruktúry namiesto ťažkého OOP.}
        Päť tried navrhovaných v~Zadaní~1 (\texttt{Station}, \texttt{TrainService},
        \texttt{StopTime}, \texttt{Leg}, \texttt{Route}) sme nahradili
        ľahkými \texttt{NamedTuple} záznamami a~predindexovanými
        \texttt{dict} mapami, aby evaluácia kandidáta zostala striktne~$O(k)$
        bez režijných nákladov na inštancovanie objektov.

  \item \textbf{Bez priebežných zastávok za behu.}
        Počas predspracovania predpočítame všetky priame platné úseky
        $(id_i,u_i,v_i,d_i,a_i)$ medzi usporiadanými dvojicami zastávok
        toho istého spoja; zoznamy medziľahlých staníc sa za behu neuchovávajú.

  \item \textbf{Viacúčelové rozšírenie ako budúca práca.}
        Pre Zadanie~2 ponechávame skalárne~$f(x)$; Pareto optimalizáciu
        samostatných kritérií~$T(x)$, $N(x)$ a~$a_k$ uvádzame v~diskusii
        budúcej práce (pozri~§4).
\end{enumerate}

%==============================================================================
\section{Implementácia modelu a~generovanie inštancií}
%==============================================================================

\subsection{Indexovaný cestovný poriadok}

Časy ukladáme ako celé minúty od $00{:}00$ dňa~$D$, čo umožňuje korektnú
aritmetiku cez polnoc. Pre každý spoj s~postupnosťou zastávok
$p_1,\ldots,p_m$ vygenerujeme všetky priame úseky $(p_a,p_b)$ s
$1\le a<b\le m$ a~$a_b>d_a$. Nad množinou úsekov budujeme dva hash indexy:
\begin{itemize}
  \item $\texttt{valid\_legs}[(D,id,u,v,d)]=a$ --- existencia a~príchod v~$O(1)$;
  \item $\texttt{departures\_from}[(D,u)]$ --- zoznam úsekov zo stanice~$u$,
        zoradený podľa~$d$, na rozšírenie susedov v~konštruktívnom hľadaní.
\end{itemize}
V~dodaných dátach má koridor + syntetika spolu 13~staníc, 16~vlakov
a~82~expandovaných úsekov (82 kľúčov v~\texttt{valid\_legs}).

\subsection{Overenie prípustnosti v~$O(k)$}

Funkcia \texttt{evaluate\_candidate} overí postupne:
\begin{enumerate}[label=\arabic*.]
  \item neprázdnosť: $k\ge 1$;
  \item koncové body: $u_1=s$, $v_k=z$;
  \item odchodové okno: $t_0\le d_1\le t_0+H$;
  \item platnosť úseku a~kladné trvanie cez \texttt{valid\_legs}, $a_i>d_i$;
  \item nadväznosť a~zmena vlaku: $v_i=u_{i+1}$, $id_i\neq id_{i+1}$;
  \item prestupný buffer: $d_{i+1}-a_i\ge\tau_{\min}$;
  \item acyklickosť: množina navštívených staníc so zamietnutím opakovaní.
\end{enumerate}
Kandidát porušujúci čo len jednu podmienku je \emph{striktne} zamietnutý
bez opravy; pri prípustnosti sa dopočíta $(T,N,f)$ a~porovnávacia trojica.

\subsection{Dátové zdroje}

\begin{itemize}
  \item \texttt{data/slovak\_corridor.json} --- koridorový podmnožina
        (priamy IC, rýchliky s~prestupom v~Žiline, nočný EN cez polnoc).
  \item \texttt{data/synthetic\_cases.json} --- riadené prípady
        (tesný prestup 4\,min, cyklus $A\to B\to A\to C$, izolovaný spoj).
  \item \texttt{scripts/parse\_gtfs.py} --- voliteľný prevodník zo surových
        GTFS CSV (\texttt{stops}/\texttt{trips}/\texttt{stop\_times}) do JSON;
        dodávka obsahuje už pripravený subset, takže grader nepotrebuje
        externé súbory.
\end{itemize}

Implementácia je v~čistom Pythone~3.12 (štandardná knižnica:
\texttt{json}, \texttt{csv}, \texttt{random}, \texttt{time}, \texttt{typing}).

%==============================================================================
\section{Overenie baseline a~experimentálne výsledky}
%==============================================================================

\subsection{Dve baseline metódy}

\begin{itemize}
  \item \textbf{Blind Random Sampling (diagnostická).}
        Náhodne vyberie $k\in\{1,2,3\}$ úsekov z~$J$ a~odošle ich do
        \texttt{evaluate\_candidate}. Dokazuje, že pri tesnej
        časovo-topologickej väzbe je miera prípustnosti blízka nule
        a~že kontrolér zachytáva všetky kategórie porušení.

  \item \textbf{Constructive Randomized Walk (primárna).}
        Vzorkuje prvý úsek zo $\texttt{departures\_from}[(D,s)]$
        s~$d_1\in[t_0,t_0+H]$; ďalej filtruje odchody z~$v_i$ podľa
        $id_{i+1}\neq id_i$, $d_{i+1}\ge a_i+\tau_{\min}$
        a~$v_{i+1}\notin\mathrm{visited}$, až kým $v_k=z$,
        nastane slepá ulička, alebo $k=k_{\max}$.
        Neúplné prechádzky sa \emph{neopravujú} --- vyhodnotia sa
        a~započítajú do štatistiky $\mathrm{feasible\_count}/\mathrm{total}$.
\end{itemize}

\subsection{Štyri verifikačné scenáre}

Všetky behy používajú \texttt{random.seed(42)}. Celý \texttt{python main.py}
dobehne za $\approx 0{,}01$\,s.

\paragraph{Scenár 1 --- trade-off penalizácie $\lambda$.}
Porovnávame pomalší priamy spoj ($T=190$, $N=0$) s~rýchlejším
jednoprestupovým spojením ($T=170$, $N=1$):
\begin{center}
\small
\begin{tabular}{@{}lcccc@{}}
\toprule
Kandidát & $T$ & $N$ & $f$ pri $\lambda{=}15$ & $f$ pri $\lambda{=}25$ \\
\midrule
A: IC521 Košice$\to$Bratislava & 190 & 0 & 190 & 190 \\
B: R801+EC125 (prestup Žilina) & 170 & 1 & \textbf{185} & 195 \\
\midrule
Preferovaný & & & B (1~prestup) & A (priamy) \\
\bottomrule
\end{tabular}
\end{center}
Pri $\lambda=15$ je výhodnejší prestup; pri $\lambda=25$ vyhráva priamy spoj ---
presne podľa $f=T+\lambda N$.

\paragraph{Scenár 2 --- prechod cez polnoc.}
Nočný EN477 Žilina$\to$Bratislava: $d_1=1395$ ($23{:}15$),
$a_k=1530$ ($01{:}30{+}1\mathrm{d}$) dáva $T=135$\,min a~prejde všetky
časové kontroly.

\paragraph{Scenár 3 --- audit neprípustných kandidátov.}
Ručne zostavené trasy porušujúce (a)~$\tau_{\min}$ (čakanie 4\,min),
(b)~odchodové okno $d_1\notin[t_0,t_0+H]$ a~(c)~acyklickosť
($A\to B\to A\to C$) sú zamietnuté s~presnými hláseniami porušení.

\paragraph{Scenár 4 --- nedosiahnuteľný cieľ.}
Neskoré okno $[1410,1440]$ z~Košíc do Bratislavy nemá žiadny prípustný
prvý odchod; konštruktívna walk hlási
\texttt{"No feasible route found"} (0/50 prípustných).

\subsection{Porovnanie baseline (Košice$\to$Bratislava, $\lambda=15$)}

\begin{center}
\small
\begin{tabular}{@{}lccc@{}}
\toprule
Metóda & Prípustné / 500 & Miera & Najlepšie~$f$ \\
\midrule
Blind Random Sampling & 3 & 0{,}6\,\% & 190 \\
Constructive Randomized Walk & 225 & 45{,}0\,\% & \textbf{185} \\
\bottomrule
\end{tabular}
\end{center}

Blind Random takmer vždy zlyhá (očakávané). Konštruktívna walk dosahuje
$\approx 45\,\%$ prípustnosť a~pri 500 rolloutoch nájde aj optimum scenára~1
($f=185$). Najčastejšie ukončenia: \texttt{feasible}~(225),
\texttt{dead\_end}~(229), \texttt{k\_max}~(46). Výsledky potvrdzujú súlad
implementácie s~matematickým modelom zo Zadania~1.

%==============================================================================
\section{Návrh metaheuristík (SA a~ACO)}
%==============================================================================

Baseline je užitočný na overenie modelu, ale náhodná konštrukcia neskúma
systematicky okolie sľubných trás. Pre Zadanie~3 navrhujeme dve
metaheuristiky, ktoré prirodzene nadväzujú na $O(k)$ evaluáciu
a~index \texttt{departures\_from}.

\subsection{Simulated Annealing (SA)}

\textbf{Zdôvodnenie.} Evaluácia kandidáta je $O(k)$ a~lacná, preto SA môže
prijať tisíce návrhov za sekundu. Deterministický greedy výber najskoršieho
odchodu často uviazne v~lokálnom optime (príliš skorý odchod s~dlhým
neskorším čakaním); SA s~teplotným plánom dokáže také riešenia dočasne
prijať a~neskôr ich vylepšiť.

\textbf{Reprezentácia a~okolia.} Stav je prípustná (alebo dočasne
mäkko-penalizovaná) postupnosť úsekov~$x$. Navrhované operátory:
\begin{itemize}
  \item \emph{subpath rerouting} --- náhodný súvislý podúsek
        $\ell_i,\ldots,\ell_j$ nahradíme novou konštruktívnou prechádzkou
        medzi $u_i$ a~$v_j$ s~respektovaním časového okna po príchode
        predchádzajúceho úseku;
  \item \emph{first-departure shift} --- prvý úsek vymeníme za iný odchod
        z~$s$ v~okne $[t_0,t_0+H]$ a~zvyšok dorátame konštruktívne;
  \item \emph{mäkká penalizácia $\tau_{\min}$} pri vysokých teplotách ---
        kandidáty s~mierne kratším prestupom ($d_{i+1}-a_i\in[0,\tau_{\min})$)
        dostanú aditívnu penalizáciu namiesto tvrdého zamietnutia; pri
        chladnutí sa tvrdý filter obnoví.
\end{itemize}
Akceptácia: $\Delta f\le 0$ vždy; inak s~pravdepodobnosťou
$e^{-\Delta f/T}$. Chladenie geometrické ($T\leftarrow \alpha T$,
$\alpha\approx 0{,}95$).

\subsection{Ant Colony Optimization (ACO)}

\textbf{Zdôvodnenie.} ACO konštruuje trasu krok za krokom rovnako ako
baseline walk, takže topologická nadväznosť ($v_i=u_{i+1}$) a~časový
buffer ($d_{i+1}-a_i\ge\tau_{\min}$) sú zaručené už pri generovaní ---
na rozdiel od slepého vzorkovania z~$J$.

\textbf{Feromóny a~heuristika.} Feromónová matica $\tau_{(u,v)}$ na
pároch staníc (prípadne jemnejšie na kľúčoch $(u,v,id)$);
viditeľnosť
\begin{equation}
  \eta_{uv}
  \propto
  \frac{1}{1+(d_{\mathrm{next}}-a_{\mathrm{curr}})+\gamma\cdot\mathbf{1}[v\neq z]},
\end{equation}
penalizuje dlhé čakanie a~odklon od cieľa. Pravdepodobnosť výberu ďalšieho
úseku je klasická
$p\propto \tau^\alpha\eta^\beta$ nad filtrovanými kandidátmi z
\texttt{departures\_from}. Elitná aktualizácia:
$\Delta\tau \propto 1/f(x^{*})$ pre najlepšiu trasu iterácie / globálne
najlepšiu, s~evaporáciou $(1-\rho)\tau$.

\paragraph{Budúca práca (viacúčelovosť).}
Po stabilizácii SA/ACO na skalárnom~$f$ plánujeme Pareto frontu
$(T,N,a_k)$ (napr.\ NSGA-II alebo ACO s~viacerými feromónovými matricami),
čo reflektuje pripomienku zo Zadania~1, že oddelenie kritérií by bolo
užitočné rozšírenie.

\vspace{0.6em}
\hrule
\vspace{0.4em}
{\small
\noindent\textbf{Príloha --- spustenie.}
\texttt{cd tim14\_zadanie2 \&\& python main.py}\\
Požiadavky: Python~3.12+, bez externých balíkov.
Zdrojové súbory: \texttt{src/model.py}, \texttt{src/instances.py},
\texttt{src/baseline.py}.
}

\end{document}
